% !TeX root = ../navier-stokes-manuscript.tex
% ============================================================
% 3.2 The Compatible Intersection
% ============================================================
\subsection{The Compatible Intersection}\label{sub:TheCompatibleIntersection}

For each finite pair \((N,M)\), the whole-terminal estimate gives one rectangle on \(I_*\).
If \(N\) or \(M\) is increased, the larger rectangle contains every derivative in the smaller one, and each shared entry is the same distributional derivative of \(u\).
Order the finite pairs coordinatewise and use the bonding maps from \eqref{eq:WholeTerminalRectangleCompatibility}.
The compatible families form
\[
  \mathscr X_\infty(I_*)
  :=
  \varprojlim_{N,M}\mathscr X_{N,M}(I_*).
\]
A family \((\mathbf R_{N,M})_{N,M}\) belongs to this set exactly when \(\mathbf R_{N,M}\in\mathscr X_{N,M}(I_*)\) for every finite pair and
\[
  \rho_{N',M';I_*}^{N,M;I_*}(\mathbf R_{N,M})
  =\mathbf R_{N',M'}
\]
whenever \(N'\leq N\) and \(M'\leq M\).

\begin{proposition}[Compatible projective assembly]\label{prop:CompatibleProjectiveAssembly}
The terminal rectangles determine the datum-specific inverse-limit element
\[
  \cI[u_0]
  :=
  \bigl(\cR_{N,M}[u_0]\bigr)_{N,M}.
\]
This family belongs to \(\mathscr X_\infty(I_*)\).
Its velocity realization is the same distribution \(u\) in every coordinate and satisfies
\begin{equation}\label{eq:SectionThreeMixedAmbientSpace}
  u\in
  \bigcap_{r=0}^{\infty}
  \bigcap_{m=0}^{\infty}
  H^r(I_*;H^m(\R^3)).
\end{equation}
Every coordinate projection commutes with every restriction map.
\end{proposition}

\begin{proof}
Fix finite \(r,m\in\Nzero\).
Choose \(N\geq r\) and \(M\geq m\).
The rectangle \(\cR_{N,M}[u_0]\) contains
\[
  \partial_t^n u\in L^2(I_*;H^{M+1})
  \hookrightarrow L^2(I_*;H^m)
  \qquad(0\leq n\leq r).
\]
These are the distributional time derivatives of the same velocity field, so \(u\in H^r(I_*;H^m)\).
The arbitrariness of \(r\) and \(m\) gives \eqref{eq:SectionThreeMixedAmbientSpace}.
Equation \eqref{eq:WholeTerminalRectangleCompatibility} makes every lower record the projection of every larger record.
The compatible family thus determines one inverse-limit element with no representative choices to reconcile.
\end{proof}

Taking the zeroth temporal projection means retaining \(u\) itself while allowing the spatial depth to range through every finite order.
Define
\begin{equation}\label{eq:SectionThreeCompleteSpatialRow}
  \Szero(I_*)
  :=\bigcap_{m=0}^{\infty}L^2(I_*;H^m(\R^3)).
\end{equation}
For each \(m\), the projection selects the velocity coordinate of \(\cR_{0,m}[u_0]\).
Compatibility says that all of these coordinates are the same \(u\), viewed in different Sobolev spaces.
Consequently,
\[
  \pi_0\bigl(\cI[u_0]\bigr)
  =u\in\Szero(I_*).
\]

This projection also contains every finite spatial energy read by the tower.
For \(\Lambda=(-\Delta)^{1/2}\) and \(s\geq0\), set
\begin{equation}\label{eq:SectionThreeSpatialEnergy}
  E_s(t):=\frac12\norm{\Lambda^su(t)}_2^2.
\end{equation}
Choosing an integer \(m\geq s\) in \eqref{eq:SectionThreeCompleteSpatialRow} gives
\begin{equation}\label{eq:SectionThreeAllSpatialEnergies}
  \int_0^{T_*}E_s(t)\,dt
  \leq
  C_{s,m}\norm{u}_{L^2(I_*;H^m)}^2
  <\infty.
\end{equation}
The intersection keeps these finite Sobolev norms as separate coordinates of the same \(u\).
